\section{The partial transpose and the PPT property}

Transposition is a positive map that is not completely positive, so applying it
to one factor of a bipartite operator can break positivity. The resulting
operation, the partial transpose
\cite[Equation~(3.16)]{Wolf2012Quantum}, sends $\rho\in\MN{d}\otimes\MN{d'}$ to
the operator $\rho^{T_1}$ obtained by transposing the first tensor factor while
leaving the second untouched, and the property $\rho^{T_1}\ge 0$---positive
partial transpose, or PPT---is the criterion
\cite[Equation~(3.17)]{Wolf2012Quantum} that separates many entangled states
from the separable ones.

\begin{definition}[Partial transpose over the first factor]
    \label{def:partial_transpose_left}
    \lean{Matrix.partialTransposeLeft}
    \leanok
    For a bipartite matrix $\rho$ on $\MN{d}\otimes\MN{d'}$, the partial transpose
    over the first factor is the matrix $\rho^{T_1}$ with entries
    \begin{align}
        (\rho^{T_1})_{(i,k),(j,l)}
        &=\rho_{(j,k),(i,l)},
        \notag
    \end{align}
    transposing the first index pair $i,j$ and fixing the second pair $k,l$.
\end{definition}

\begin{definition}[Partial transpose over the second factor]
    \label{def:partial_transpose_right}
    \lean{Matrix.partialTransposeRight}
    \leanok
    For a bipartite matrix $\rho$ on $\MN{d}\otimes\MN{d'}$, the partial transpose
    over the second factor is the matrix $\rho^{T_2}$ with entries
    \begin{align}
        (\rho^{T_2})_{(i,k),(j,l)}
        &=\rho_{(i,l),(j,k)},
        \notag
    \end{align}
    transposing the second index pair $k,l$ and fixing the first pair $i,j$.
\end{definition}

\begin{theorem}[The partial transpose is an involution]
    \label{thm:partial_transpose_left_involution}
    \lean{Matrix.partialTransposeLeft_partialTransposeLeft}
    \leanok
    \uses{def:partial_transpose_left}
    Applying the first-factor partial transpose twice returns the original matrix:
    $(\rho^{T_1})^{T_1}=\rho$.
\end{theorem}

\begin{proof}
    \leanok
    Both entries equal $\rho_{(i,k),(j,l)}$: the first transposition exchanges the
    first index pair, the second restores it.
\end{proof}

\begin{theorem}[The partial transpose preserves the trace]
    \label{thm:trace_partial_transpose_left}
    \lean{Matrix.trace_partialTransposeLeft}
    \leanok
    \uses{def:partial_transpose_left}
    The first-factor partial transpose preserves the trace:
    $\tr(\rho^{T_1})=\tr(\rho)$.
\end{theorem}

\begin{proof}
    \leanok
    On the diagonal the row and column indices coincide, so the index exchange in
    the definition is the identity and each diagonal entry is unchanged.
\end{proof}

\begin{theorem}[The partial transpose preserves Hermiticity]
    \label{thm:partial_transpose_left_hermitian}
    \lean{Matrix.IsHermitian.partialTransposeLeft}
    \leanok
    \uses{def:partial_transpose_left}
    If $\rho$ is Hermitian then so is its first-factor partial transpose
    $\rho^{T_1}$.
\end{theorem}

\begin{proof}
    \leanok
    The entry $(\rho^{T_1})_{(i,k),(j,l)}=\rho_{(j,k),(i,l)}$ has complex
    conjugate $\overline{\rho_{(j,k),(i,l)}}=\rho_{(i,l),(j,k)}$ by Hermiticity of
    $\rho$, which is the $((j,l),(i,k))$ entry of $\rho^{T_1}$.
\end{proof}

\begin{theorem}[The two partial transposes differ by a global transposition]
    \label{thm:partial_transpose_right_eq_transpose_left}
    \lean{Matrix.partialTransposeRight_eq_transpose_partialTransposeLeft}
    \leanok
    \uses{def:partial_transpose_left, def:partial_transpose_right}
    The second-factor partial transpose is the full transpose of the first-factor
    one: $\rho^{T_2}=(\rho^{T_1})^{\mathsf T}$.
\end{theorem}

\begin{proof}
    \leanok
    Transposing both factors is the full transpose, so transposing the first
    factor and then the whole operator transposes only the second factor.
\end{proof}

\begin{definition}[Positive partial transpose property]
    \label{def:is_ppt}
    \lean{Matrix.IsPPT}
    \leanok
    \uses{def:partial_transpose_left}
    A bipartite matrix $\rho$ has the positive-partial-transpose (PPT) property when
    its first-factor partial transpose is positive semidefinite, $\rho^{T_1}\ge 0$.
\end{definition}

\begin{theorem}[The PPT property is symmetric in the two factors]
    \label{thm:is_ppt_iff_right}
    \lean{Matrix.isPPT_iff_partialTransposeRight_posSemidef}
    \leanok
    \uses{def:is_ppt, def:partial_transpose_right,
        thm:partial_transpose_right_eq_transpose_left}
    A bipartite matrix has the PPT property if and only if its second-factor
    partial transpose is positive semidefinite:
    $\rho^{T_1}\ge 0\iff\rho^{T_2}\ge 0$.
\end{theorem}

\begin{proof}
    \leanok
    The two partial transposes differ by a global transposition, and transposition
    preserves positive semidefiniteness, so one is positive exactly when the other
    is.
\end{proof}

\begin{theorem}[Basis-change law for the partial transpose]
    \label{thm:partial_transpose_left_basis_change}
    \lean{Matrix.partialTransposeLeft_basis_change}
    \leanok
    \uses{def:partial_transpose_left}
    For matrices $U,V$ on the first factor and a bipartite matrix $\rho$,
    conjugating inside by $V$ and $U$, taking the first-factor partial transpose,
    and conjugating outside by $U$ and $V$ produces the original partial transpose
    conjugated by $U U^{\mathsf T}$ and $V^{\mathsf T}V$:
    \begin{align}
        (U\otimes\Id)
        \left[((V\otimes\Id) \rho\,(U\otimes\Id))^{T_1}\right]
        (V\otimes\Id)
        &=((U U^{\mathsf T})\otimes\Id) \rho^{T_1}
            ((V^{\mathsf T}V)\otimes\Id).
        \notag
    \end{align}
\end{theorem}

\begin{proof}
    \leanok
    Expanding the conjugated operator entry by entry as a double sum over the first
    factor, the partial transpose exchanges the two summation indices and reverses
    the roles of the left and right conjugating matrices, leaving them transposed;
    collecting the Kronecker factors gives the stated form.
\end{proof}

\begin{theorem}[Conjugation invariance of the partial transpose]
    \label{thm:is_ppt_basis_change}
    \lean{Matrix.isPPT_basis_change}
    \leanok
    \uses{def:is_ppt, thm:partial_transpose_left_basis_change}
    Let $U$ be any matrix on the first factor. If $\rho$ has the PPT property then
    \begin{align}
        &(U\otimes\Id)
        \Bigl[\bigl((U^\dagger\otimes\Id) \rho
        (U\otimes\Id)\bigr)^{T_1}\Bigr]
        (U^\dagger\otimes\Id)
        \notag
    \end{align}
    is positive semidefinite; no invertibility or unitarity of $U$ is required.
    Specializing to a unitary $U$ this is the partial transpose taken in the basis
    changed by $U$, so the positivity of the partial transpose is independent of the
    local basis, recovering Wolf's basis-independence statement
    \cite[Equation~(3.17)]{Wolf2012Quantum}.
\end{theorem}

\begin{proof}
    \leanok
    By the basis-change law with $V=U^\dagger$, the new-basis partial transpose
    equals $((U U^{\mathsf T})\otimes\Id) \rho^{T_1}
    ((U U^{\mathsf T})\otimes\Id)^\dagger$, a conjugation of $\rho^{T_1}$
    that preserves positive semidefiniteness.
\end{proof}

\begin{definition}[Transposition witness]
    \label{def:transposition_witness}
    \lean{Matrix.transpositionWitness}
    \leanok
    For matrices $A,B$ on $\MN{d}$ and the SWAP operator $F$ on
    $\MN{d}\otimes\MN{d}$, the transposition witness is
    $W=(A\otimes B)\,F\,(A\otimes B)^\dagger$.
    These are the entanglement witnesses associated with the transposition map.
\end{definition}

\begin{theorem}[Transposition witnesses are Hermitian]
    \label{thm:transposition_witness_hermitian}
    \lean{Matrix.transpositionWitness_isHermitian}
    \leanok
    \uses{def:transposition_witness}
    Every transposition witness $W=(A\otimes B)\,F\,(A\otimes B)^\dagger$ is
    Hermitian.
\end{theorem}

\begin{proof}
    \leanok
    The SWAP operator is self-adjoint, $F^\dagger=F$, so taking the adjoint of
    $(A\otimes B)\,F\,(A\otimes B)^\dagger$ returns the same operator.
\end{proof}

\section{Separable states and the PPT criterion}

A bipartite state is separable when it is a convex combination of pure product
states $\sum_i p_i\,\ket{a_i}\bra{a_i}\otimes\ket{b_i}\bra{b_i}$. Each pure product
state is a Kronecker product of two rank-one positive semidefinite matrices, the
non-negative weight $p_i$ is absorbed into one factor, and conversely a positive
semidefinite matrix is a non-negative combination of rank-one projectors. Up to
normalization this gives the equivalent description of separable states as finite
sums of Kronecker products of positive semidefinite matrices, the cone generated by
the product positive operators, which is the form used below.

\begin{definition}[Separable bipartite matrix]
    \label{def:is_separable}
    \lean{Matrix.IsSeparable}
    \leanok
    A bipartite matrix $\rho$ on $\MN{d}\otimes\MN{d'}$ is separable when it is a
    finite sum of Kronecker products of positive semidefinite matrices,
    \begin{align}
        \rho&=\sum_i A_i\otimes B_i,
        \notag\\
        A_i&\ge 0,
        \notag\\
        B_i&\ge 0.
        \notag
    \end{align}
\end{definition}

\begin{theorem}[A Kronecker product of positive operators is separable]
    \label{thm:is_separable_kronecker}
    \lean{Matrix.isSeparable_kronecker}
    \leanok
    \uses{def:is_separable}
    For positive semidefinite $A$ on $\MN{d}$ and $B$ on $\MN{d'}$, the Kronecker
    product $A\otimes B$ is separable.
\end{theorem}

\begin{proof}
    \leanok
    It is the one-term sum $\sum_i A_i\otimes B_i$ with the single index taking the
    value $A\otimes B$.
\end{proof}

\begin{theorem}[Separable matrices are closed under addition]
    \label{thm:is_separable_add}
    \lean{Matrix.IsSeparable.add}
    \leanok
    \uses{def:is_separable}
    If $\rho$ and $\sigma$ are separable then so is $\rho+\sigma$.
\end{theorem}

\begin{proof}
    \leanok
    Concatenating the two index families realizes $\rho+\sigma$ as a single sum of
    Kronecker products of positive semidefinite matrices.
\end{proof}

\begin{theorem}[Separable matrices are closed under non-negative scaling]
    \label{thm:is_separable_smul}
    \lean{Matrix.IsSeparable.smul}
    \leanok
    \uses{def:is_separable}
    If $\rho$ is separable and $c\ge 0$ is real then $c\,\rho$ is separable.
\end{theorem}

\begin{proof}
    \leanok
    Absorbing $c$ into the first factor of each Kronecker product keeps every factor
    positive semidefinite, since a non-negative multiple of a positive semidefinite
    matrix is positive semidefinite.
\end{proof}

\begin{theorem}[A separable matrix is positive semidefinite]
    \label{thm:is_separable_pos_semidef}
    \lean{Matrix.IsSeparable.posSemidef}
    \leanok
    \uses{def:is_separable}
    Every separable matrix is positive semidefinite.
\end{theorem}

\begin{proof}
    \leanok
    Each summand $A_i\otimes B_i$ is a Kronecker product of positive semidefinite
    matrices, hence positive semidefinite, and a finite sum of positive semidefinite
    matrices is positive semidefinite.
\end{proof}

\begin{theorem}[The partial transpose distributes over a finite sum]
    \label{thm:partial_transpose_left_sum}
    \lean{Matrix.partialTransposeLeft_sum}
    \leanok
    \uses{def:partial_transpose_left}
    For a finite family of bipartite matrices,
    $(\sum_i f_i)^{T_1}=\sum_i f_i^{T_1}$.
\end{theorem}

\begin{proof}
    \leanok
    The first-factor partial transpose is an entrywise reindexing of the matrix, so
    it commutes with the entrywise finite sum.
\end{proof}

\begin{theorem}[PPT criterion, forward direction]
    \label{thm:is_separable_is_ppt}
    \lean{Matrix.IsSeparable.isPPT}
    \leanok
    \uses{def:is_separable, def:is_ppt, thm:partial_transpose_left_sum}
    Every separable matrix has the positive-partial-transpose property:
    if $\rho$ is separable then $\rho^{T_1}\ge 0$. A negative partial transpose
    therefore certifies that a state is entangled.
\end{theorem}

\begin{proof}
    \leanok
    The first-factor partial transpose is linear and acts on a Kronecker product by
    transposing the first factor, $(A\otimes B)^{T_1}=A^{\mathsf T}\otimes B$.
    Writing $\rho=\sum_i A_i\otimes B_i$ with each factor positive semidefinite, the
    partial transpose is $\sum_i A_i^{\mathsf T}\otimes B_i$. The transpose of a
    positive semidefinite matrix is positive semidefinite, so each summand is a
    Kronecker product of positive semidefinite matrices, and the sum is positive
    semidefinite.
\end{proof}

\begin{theorem}[Applying a map to one factor distributes over a finite sum]
    \label{thm:tensor_map_id_sum}
    \lean{Matrix.tensorMapId_sum}
    \leanok
    \uses{def:tensor_map_id}
    For a linear map $T$ on the first factor and a finite family of bipartite
    matrices,
    $(T\otimes\id)(\sum_i f_i)=\sum_i (T\otimes\id)(f_i)$.
\end{theorem}

\begin{proof}
    \leanok
    Applying a map to the first factor is linear in the bipartite matrix, so it
    commutes with the finite sum.
\end{proof}

\begin{theorem}[Reduction map at the first level on a separable state]
    \label{thm:is_separable_reduction_one}
    \lean{Matrix.IsSeparable.tensorMapId_tEta_one_posSemidef}
    \leanok
    \uses{def:is_separable, def:t_eta, thm:reduction_map_one_positive,
        thm:tensor_map_id_sum}
    For $d\ge 1$ and a separable matrix $\rho$ on $\MN{d}\otimes\MN{d'}$, applying
    the reduction map $T_1(X)=\tr(X)\Id-X$ to the first factor keeps the result
    positive semidefinite: $(T_1\otimes\id)(\rho)\ge 0$.
    This is the separability instance of the operator inequality underlying the
    reduction criterion at the first level.
\end{theorem}

\begin{proof}
    \leanok
    The reduction map $T_1$ is positive, so it maps each positive semidefinite first
    factor $A_i$ of $\rho=\sum_i A_i\otimes B_i$ to a positive semidefinite matrix.
    The ampliation acts by $(T_1\otimes\id)(A_i\otimes B_i)=T_1(A_i)\otimes B_i$, each
    summand is a Kronecker product of positive semidefinite matrices, and the sum is
    positive semidefinite.
\end{proof}
